\begin{table}[htb!]
\centering
\footnotesize
\begin{tabular}{@{}lrrrrrrrr@{}}
\toprule
\textbf{Pair} & \textbf{9,900} & \textbf{10,000} & \textbf{10,050} & \textbf{10,100} & \textbf{10,200} & \textbf{Mean} & \textbf{SD} & \textbf{Relative SD} \\
\midrule
AUD/USD & 94.2\% & 87.2\% & 117.1\% & 108.3\% & 105.7\% & 102.5\% & 10.58 & 10.3\% \\
EUR/USD & 69.5\% & 70.6\% & 76.5\% & 77.8\% & 75.1\% & 73.9\% & 3.27 & 4.4\% \\
GBP/USD & 126.5\% & 129.1\% & 127.9\% & 122.0\% & 115.1\% & 124.1\% & 5.11 & \textbf{4.1\%} \\
NZD/USD & 61.7\% & 60.0\% & 54.4\% & 53.7\% & 52.1\% & 56.4\% & 3.76 & 6.7\% \\
USD/CAD & 230.4\% & 221.0\% & 218.5\% & 216.9\% & 246.0\% & 226.6\% & 10.81 & 4.8\% \\
USD/CHF & 56.4\% & 54.8\% & 43.4\% & 52.6\% & 45.5\% & 50.6\% & 5.16 & 10.2\% \\
USD/JPY & 15.9\% & 16.2\% & 16.4\% & 15.7\% & 14.5\% & 15.7\% & 0.69 & 4.4\% \\
\bottomrule
\end{tabular}
\caption{Total return from each of the five starting balances, in euros. All thirty-five runs are profitable. Because broker volumes are discrete, a different opening balance rounds positions differently and produces a genuinely different sequence of trades rather than a rescaled copy of the same one.}
\label{Tables:PathRobustness}
\end{table}